\documentclass[twocolumn,superscriptaddress,unsortedaddress,nofootinbib,amsmath,amssymb,
aps,floatfix]{revtex4-2}
\usepackage{graphicx,tikz} 
\usetikzlibrary{arrows.meta,decorations.pathreplacing,calc,positioning}
\usepackage[mode=buildmissing]{standalone} 
\usepackage{dcolumn}
\usepackage{bm}
\usepackage{physics}
\usepackage{xcolor}
\usepackage{siunitx}
\usepackage{soul}
\usepackage{comment}
\usepackage[bb=boondox]{mathalfa}
\usetikzlibrary{external}
\tikzexternalize[prefix=figures/]
\DeclareMathOperator{\sinc}{sinc}
\newcommand{\may}[1]{\textcolor{orange}{#1}}
\begin{document}
\title{Lipkin model, Entanglement, Noise, Dephasing and dissipation}

\author{Juan David Álvarez-Cuartas}
\affiliation{Departament of Physics and Center for Bioinformatics and Photonics---CIBioFi, 
Universidad del Valle, Cali 760032, Colombia}
\affiliation{Department of Physics and Center for Computing in Science Education, University of Oslo, N-0316 Oslo, Norway}
\author{Patrick Cook and Danny Jammooa}
\affiliation{Facility for Rare Isotope Beams and Department of Physics and Astronomy, Michigan State University, East Lansing, Michigan 48824, USA}
\author{Morten Hjorth-Jensen}
\affiliation{Department of Physics and Center for Computing in Science Education, University of Oslo, N-0316 Oslo, Norway}
\author{Dean Lee}
\affiliation{Facility for Rare Isotope Beams and Department of Physics and Astronomy, Michigan State University, East Lansing, Michigan 48824, USA}
\author{John H.~Reina, and Juan M.~Scarpetta}
\affiliation{Departament of Physics and Center for Bioinformatics and Photonics---CIBioFi, 
Universidad del Valle, Cali 760032, Colombia}
\affiliation{Department of Physics and Center for Computing in Science Education, University of Oslo, N-0316 Oslo, Norway}
\begin{abstract}
Text to come here
\end{abstract}

%\pacs{02.70.Ss, 31.15.A-, 31.15.bw, 71.15.-m, 73.21.La}


\date{\today}

\maketitle

\section{Introduction} 
Structure of article
\begin{enumerate}
    \item Introduction with motivation
    \item Discussion of the Lipkin model
    \item Setting up time evolution
    \item PMMs and time evolution
    \item 
\end{enumerate}

\section{Mathematical Framework and Entanglement Measures}

\subsection{Hamiltonian and Symmetry.}
The Lipkin--Meshkov--Glick (LMG) model describes $N$ spin-$\tfrac{1}{2}$ particles interacting through collective spin operators.  
Its Hamiltonian can be expressed as
\begin{equation}
  \hat{H}_{\mathrm{LMG}} = -\frac{\lambda}{N} \left( \hat{S}_x^2 + \gamma\, \hat{S}_y^2 \right)
  - h\, \hat{S}_z,
  \label{eq:lmg-hamiltonian}
\end{equation}
where $\hat{S}_\alpha = \tfrac{1}{2}\sum_{i=1}^N \hat{\sigma}_i^\alpha$ are collective spin operators,
$\lambda$ is the interaction strength, $\gamma$ the anisotropy parameter, and $h$ an external magnetic field.
The factor $1/N$ ensures a well-defined thermodynamic limit.
The total spin $\hat{\mathbf{S}}^2$ commutes with the Hamiltonian, implying conservation of total spin $S = N/2$,
and the system therefore lives in the symmetric subspace spanned by the Dicke states
\begin{equation}
  \ket{S, M} \equiv \ket{N/2, M}, \quad M = -S, \ldots, S.
\end{equation}


%\documentclass[tikz,border=6pt]{standalone}

%\begin{figure}
\begin{tikzpicture}[font=\small, every node/.style={inner sep=1pt}]
  % Levels
  \coordinate (Lleft) at (0,0);
  \coordinate (Rright) at (8,0);
  \draw[thick] (0,2) -- (6,2) node[right]{\(\text{upper level}\,(\,+\tfrac{\varepsilon}{2}\,)\)};
  \draw[thick] (0,0) -- (6,0) node[right]{\(\text{lower level}\,(\,-\tfrac{\varepsilon}{2}\,)\)};

  % Particles in lower level: four filled circles
  \foreach \i in {0,1,2,3}{
    \node[circle,fill=black,minimum size=6pt,inner sep=0pt] (p\i) at (0.8+ \i*1.0 , -0.0) {};
  }
  % Label occupancy
  \node[anchor=west] at (6.2,0) {initial: \(n_\text{up}=0,\; n_\text{low}=4\)};
  
  % Pair-scattering arrows: two-particle excitation to upper level
  \draw[->, >=Latex, thick, blue] ($(p0)+(0,0.15)$) .. controls (1.2,0.9) and (2.0,1.6) .. ($(p0)+(1.2,2)$)
     node[midway, sloped, above]{\(\tfrac{V}{2}\) pair flip};
  \draw[->, >=Latex, thick, blue] ($(p1)+(0,0.15)$) .. controls (1.6,0.9) and (2.8,1.6) .. ($(p1)+(2.6,2)$);
  
  % Draw two particles in upper level as a possible excited configuration
  \node[circle,fill=black,minimum size=6pt,inner sep=0pt] (u0) at (3.2,2.0) {};
  \node[circle,fill=black,minimum size=6pt,inner sep=0pt] (u1) at (4.2,2.0) {};
  \node[anchor=west] at (6.2,2) {excited: \(n_\text{up}=2,\; n_\text{low}=2\)};
  
  % Curved double arrow indicating pair scattering V <-> (both directions)
  \draw[<->, >=Latex, very thick, red] (3.6,2.2) .. controls (4.8,1.0) and (4.8,-0.5) .. (3.6,-0.2)
     node[midway, right, xshift=4pt] {\(\)};

  % Label the scattering more explicitly
  \node[red,right=2pt of $(3.6,0.6)$] {\(\displaystyle V\;:\; J_+^2+J_-^2\)};
  
  % Pseudo-spin mapping box
  \begin{scope}[shift={(0,-2.5)}]
    \node[draw,rounded corners, inner sep=6pt] (spinbox) at (3,0) {
      \begin{minipage}{7.5cm}
        \centering
        {\bf Pseudo-spin representation (collective)}\\[4pt]
        \(
          J_z=\tfrac{1}{2}\,(n_\text{up}-n_\text{low}),\qquad j=\tfrac{N}{2}=2
        \)\\[4pt]
        Basis: \(|j,m\rangle\) with \(m=-2,-1,0,1,2\) (e.g. \(|2,-2\rangle\) = all particles in lower level)
      \end{minipage}
    };
  \end{scope}

  % Hamiltonian near the diagram
  \node[anchor=west] at (6.6, -1.5) {
    \begin{minipage}{4.0cm}
      \[
        H = \varepsilon J_z + \frac{V}{2}\left(J_+^2 + J_-^2\right)
      \]
      \vspace{4pt}
      Notes: pairwise scattering moves two particles simultaneously between levels.
    \end{minipage}
  };

  % Decorative: arrows linking the spinbox to levels
  \draw[->, dashed] (spinbox.north) .. controls (3.0,0.3) .. (2.2,0.8) node[midway, left] {collective map};

  % Title
  \node[above left=0.5cm and -0.2cm of current bounding box.north east, anchor=north east, font=\bfseries\large] (title) {Lipkin model — $N=4$ (schematic)};
\end{tikzpicture}
%\end{figure}


\subsection{Mean-Field and Quantum Phase Transition.}
In the thermodynamic limit ($N \to \infty$), a mean-field treatment becomes exact.
A semiclassical energy functional can be written in terms of collective spin variables
$(\theta, \phi)$ on the Bloch sphere as
\begin{equation}
  E(\theta, \phi) = -\frac{\lambda}{2}\sin^2\theta \left( \cos^2\phi + \gamma\, \sin^2\phi \right)
  - h\, \cos\theta.
\end{equation}
The ground state undergoes a second-order quantum phase transition at $h_c = \lambda(1+\gamma)/2$,
separating a symmetric phase ($\langle S_x \rangle = 0$) from a symmetry-broken phase ($\langle S_x \rangle \neq 0$).
Near this critical point, the entanglement entropy exhibits logarithmic scaling with system size, signalling a diverging correlation length in the collective degrees of freedom.

\subsection{Reduced Density Matrices and Entropy.}
Because of full permutation symmetry, the reduced density matrix of any single spin
is identical and can be obtained from the collective spin density operator $\rho$ as
\begin{equation}
  \rho_1 = \mathrm{Tr}_{2,\ldots,N}(\rho)
  = \frac{1}{2}\left( \mathbb{I} + \frac{2}{N}\langle \hat{\mathbf{S}} \rangle \cdot \boldsymbol{\sigma} \right).
\end{equation}
The single-spin entanglement (or \emph{one-tangle}) with the remainder of the system is quantified by the linear entropy
\begin{equation}
  \tau_1 = 4 \det(\rho_1) = 1 - \frac{4}{N^2} \langle \hat{\mathbf{S}} \rangle^2.
\end{equation}
This quantity vanishes in product states and reaches unity for maximally mixed single-spin states.

For bipartite entanglement between two subsystems of sizes $L$ and $N-L$, the von Neumann entropy
\begin{equation}
  S_L = - \mathrm{Tr}\, (\rho_L \log_2 \rho_L)
\end{equation}
is the standard measure.  In the symmetric Dicke basis, $\rho_L$ is block-diagonal and
its eigenvalues can be expressed analytically in terms of Clebsch--Gordan coefficients,
allowing explicit evaluation of $S_L$ for moderate $N$.

\subsection{Two-Spin Correlations and Concurrence.}
Pairwise entanglement between two spins is characterized by the \emph{concurrence} $C$,
defined from the two-spin reduced density matrix $\rho_{12}$ as
\begin{equation}
  C = \max \left\{ 0, \lambda_1 - \lambda_2 - \lambda_3 - \lambda_4 \right\},
\end{equation}
where $\lambda_i$ are the square roots of the eigenvalues (in descending order)
of the matrix $\rho_{12} \tilde{\rho}_{12}$, with
$\tilde{\rho}_{12} = (\sigma_y \otimes \sigma_y)\, \rho_{12}^*\, (\sigma_y \otimes \sigma_y)$.
In the LMG model, due to full permutation symmetry, $C$ is identical for all spin pairs.
It reaches a maximum near the critical point and decays as $N^{-1/3}$ at criticality,
indicating that bipartite entanglement is sub-extensive even though the total
multipartite entanglement diverges.

\subsection{Multipartite and Mixed-State Measures.}
Beyond bipartite partitions, the \emph{global entanglement} and the
\emph{geometric measure of entanglement} are often used:
\begin{equation}
  E_{\mathrm{G}} = 1 - \max_{\ket{\phi_{\mathrm{sep}}}} |\langle \phi_{\mathrm{sep}} | \psi_0 \rangle|^2,
\end{equation}
where $\ket{\phi_{\mathrm{sep}}}$ is the closest separable state to the ground state $\ket{\psi_0}$.
For mixed or non-complementary partitions, the logarithmic negativity
\begin{equation}
  \mathcal{N} = \log_2 \| \rho^{T_A} \|_1
\end{equation}
quantifies entanglement by means of the trace norm of the partially transposed density matrix $\rho^{T_A}$.

\subsection{Scaling and Universality.}
Near the quantum critical point, different entanglement measures show universal scaling with system size $N$:
\begin{equation}
  S_L \sim \frac{1}{6} \log_2 N + \mathrm{const}, \qquad
  \tau_1 \sim N^{-2/3}, \qquad
  C \sim N^{-1/3}.
\end{equation}
These exponents coincide with those predicted by conformal field theory
for one-dimensional critical systems, even though the LMG model is infinite-range.
Such universality underscores that entanglement, rather than spatial locality,
captures the essential signatures of quantum criticality.

\section{Results and discussions}


\section{Conclusion}

%\bibliography{paper}
\end{document}

